% SOURCE docs/release/2026-10-29/source/docs/writing-repair/2026-10-29/src/bd9020f20dab21be.tex CHAR 0
\chapter{What a shared service must preserve}
\label{r3:software}

The preceding theory makes software selection more demanding than choosing
an estimator with an attractive name. A deposit-retention analysis may ask
for an average effect across customers, across branches, or across balances.
A collections analysis may ask for an effect of a fixed schedule, a rule
using available history, or a modification of the historical assignment
law. These distinctions change the estimand before any package is called.
The common product needs to preserve them as it moves between data engines,
identification libraries, statistical estimators and decision solvers.
Chapter~\ref{scisyn:software} gives the principal selection argument. Here the
comparison develops the conditions behind that choice: the population carried
by weights, the intervention represented by a policy, the information available
at each decision and the sampling unit carried into uncertainty. These are the
properties that the operations in Chapter~\ref{r3:product} must preserve.

This chapter uses the preserved primary documentation and selected source
interfaces of 44 projects. The scientific comparisons concern the inspected
versions, using 7 October 2026 as the publication/version cutoff where a
dated revision could be established. Mutable documentation accessed on
8 October has a separate access identity and does not establish its exact
contents at the cutoff. A documented interface, a successful
package execution and a validated banking analysis are different evidence.
The broad comparisons establish source/document evidence. Selected installed
DuckDB and Polars revision operations, scikit-learn Lasso normalization and
SciPy optimization now have executed reference comparisons, recorded
separately from the historical documentation snapshots. Their versions and
local source/license copies accompany the receipts; qualification applies
to those operations, not every feature of each package
\citep{runtime_duckdb,runtime_polars,runtime_scikit_learn,runtime_scipy,runtime_numpy}.
The selected DoubleML score and Tigramite temporal-search executions in
Sections~\ref{s5:composition} and~\ref{s5:discovery} add substantive causal
operations with matched references and preserved endpoint semantics
\citep{runtime_doubleml_scope,runtime_tigramite_scope}.
Published package versions, mutable documentation snapshots and
commit-pinned source files are kept distinct in the accompanying source
register. No speed ranking is inferred from repository activity or a
maintainer's benchmark.

\section{Changing software must preserve the economic quantity}
\label{r3:software-types}

Let a request contain a causal model class $\mathcal M$, observation operator
$K$, intervention $g$, target population $\nu$ and functional
$\theta(M;g,\nu)$. Let $P=K(M)$ denote the observed law. The first possible
output is an identification statement
\[
 \theta(M;g,\nu)=\Psi(K(M))\qquad\text{for all }M\in\mathcal M.
\]
An estimator is then a numerical map from the observed sample to
$\widehat\Psi$. An inference method is a further construction whose
validity depends on the sampling scheme and estimator remainder.
The input of a planner is different again: it requires an interventional
transition/observation/reward model, or an explicitly defined ambiguity set.
None of these objects can be substituted merely because its serialized
form is an array of numbers.

DoWhy exposes this separation in its interface. The inspected
\texttt{identify\_effect} selects its ID implementation when requested;
\texttt{estimate\_effect} separately takes an identified expression and
an estimator method \citep{sw_dowhy}. An expression produced by graphical
identification need not have a matching numerical estimator in the chosen
backend. Thus the product must distinguish three outcomes: a supported
identified functional, a valid identification result whose estimation is
not implemented, and nonidentification in the assumed model class.
The last requires a scientific witness such as the paired models of
Section~\ref{r3:nonidentification}; a missing Python method is not such a
witness.

There is a useful criterion for a data adapter. Suppose it transforms the
observation through $T$ and the backend estimates $\Psi_B$. Target
preservation requires
\begin{equation}
 \Psi(P)=\Psi_B(T_\#P)
 \quad\text{for every allowed observed law }P,
 \label{r3:adapter-identity}
\end{equation}
where $T_\#P$ is the distribution after transformation. Invertible relabeling
usually satisfies this with an appropriate inverse interpretation.
Demeaning, history truncation, averaging time bins, hard record linkage,
weight normalization and imputation generally require additional arguments.
Equation~\eqref{r3:adapter-identity} also shows why a common
\texttt{fit/predict} signature is insufficient: it says nothing about
which functional survives the conversion.

\section{Effects, projections and the population carried by weights}
\label{r3:software-effects}

EconML, DoubleML, causalml and grf provide useful but overlapping learning
components \citep{sw_econml,sw_doubleml,sw_causalml,sw_grf}. The relevant
comparison begins with the target. For binary treatment under adjustment,
the CATE is $\tau(x)=\E[Y(1)-Y(0)\mid X=x]$. A partially linear residual
coefficient instead solves an unconditional moment and, with heterogeneous
effects, can weight $\tau(x)$ by residual treatment variation. The
R-learner calculation in Section~\ref{r3:heterogeneity} identifies that
difference. An IV learner introduces yet another moment and its exclusion,
relevance and effect-heterogeneity restrictions. A common effect-prediction
interface cannot turn these coefficients into a common CATE.

The inspected \texttt{grf::causal\_forest} documents the distinction
explicitly: for binary unconfounded treatment it targets a CATE; for
continuous treatment it uses
\[
 b(x)=\frac{\operatorname{Cov}(Y,A\mid X=x)}
                 {\operatorname{Var}(A\mid X=x)}.
\]
This is a local linear projection in the treatment dimension. It is not
generally the dose-response derivative at a nominated action. For example,
if $A$ is uniform on $[-1,1]$, $Y(a)=a^3$, and no confounding is present,
then $b=\E[A^4]/\E[A^2]=3/5$, while the derivative at $a=0$ is zero.
Both quantities can be meaningful; an interface must name the requested one.

Population weighting is equally substantive. If branch $c$ contains $n_c$
customers with average effect $\tau_c$, then
\[
 \theta_{\rm customer}=
 \frac{\sum_c n_c\tau_c}{\sum_c n_c},
 \qquad
 \theta_{\rm branch}=\frac1C\sum_c\tau_c.
\]
For two branches with sizes $1$ and $9$ and effects $0$ and $10$, these
values are $9$ and $5$. The source of \texttt{causal\_forest} states that
\texttt{equalize.cluster.weights=TRUE} uses inverse cluster-size weights
when estimating average effects; it also requires
\texttt{sample.weights=NULL}. The default gives equal weight to units.
This option changes the population aggregation, not merely the reported
standard error \citep{sw_grf}.

A balance-weighted bank target needs a further definition. A weight
measured before treatment can define a baseline exposure-weighted average.
A balance measured after treatment may itself change under intervention;
using it as an ordinary sampling weight can select on a treatment
descendant. The product must carry the weight's time, role and normalization
with the target. Frequency weights, sampling weights, treatment-density
ratios, loss weights and numerical importance weights cannot share a
single undocumented field.

DoubleML's inspected data classes distinguish outcomes, treatments,
instruments, covariates and cluster columns. They do not make arbitrary
dependent records satisfy a cross-fitting theorem. causalml's documentation
supplies a CATE/uplift and benchmarking route; using its predictions for
contact allocation still requires the policy-value argument and independent
evaluation of Chapter~\ref{r3:decision-learning}. A package comparison should
therefore hold the estimand, population, nuisance information and evaluation
sample fixed before comparing bias, coverage, policy value or runtime.

\section{A panel identifier does not specify a panel model}
\label{r3:software-panels}

The same customer-period table can support a fixed-effect projection,
a staggered-adoption ATT, a dynamic GMM coefficient or a sequential blip.
Chapter~\ref{r3:panels} derives different assumptions for each. The software
mirrors this division. linearmodels supplies linear panel, IV and system
baselines; PyFixest and fixest supply formula-based fixed-effect, IV and
other regression operations; pydynpd explicitly constructs difference and
system GMM \citep{sw_linearmodels,sw_pyfixest,sw_fixest,sw_pydynpd}.
Absorbing many fixed effects quickly does not repair an invalid
heterogeneous-treatment TWFE estimand. Likewise, an automated search across
instrument lags cannot validate the orthogonality of those instruments.

For a branch rollout, \texttt{did::att\_gt} represents first-adoption
cohorts and group--time effects. Its inspected interface has important
implications for a shared service. Time-varying sampling weights are
handled differently in balanced panels and repeated cross sections or
unbalanced panels. Fixing weights at an anchor period can remove units
unobserved at that anchor. Beyond individual clustering, this snapshot
supports at most one additional coarser cluster variable. Uniform bands
require its bootstrap; setting \texttt{compute\_inffunc=FALSE} produces
point estimates without inference and prevents the documented aggregation
through \texttt{aggte} \citep{sw_did}. Saving memory by dropping influence
information therefore also removes downstream operations.

DoubleML's panel container is not a substitute for an arbitrary sequential
history object. Its default \texttt{static\_panel=False} interprets
treatment values as first-treatment timing, and the inspected constructor
permits only one treatment column \citep{sw_doubleml}. An array of monthly
credit-limit changes is not a cohort-date array. The conversion must be
rejected unless the requested analysis explicitly defines an appropriate
adoption representation.

DynamicDML illustrates the separate chronological contract. The inspected
EconML implementation requires equal group lengths and interprets each
group's occurrence order as its period order. Only first-period $X$
defines effect heterogeneity; later $X$ is a time-varying control. The
fit method warns that \texttt{sample\_weight} and \texttt{sample\_var}
are ignored and passes neither into fitting \citep{sw_econml}. A bank
adapter must therefore sort by a validated decision clock, check for
missing or duplicated periods, use a group-aware split, and reject a
weighted request this interface cannot honor. Padding a missing decision
with zeros changes the observation law and is not a general repair for
an unbalanced panel.

Alternatives change the identifying model rather than merely the optimizer.
synthdid's documented input requires simultaneous treatment with a block
adoption matrix. gsynth imputes untreated outcomes through interactive
fixed effects; its current README says that it wraps \texttt{fect} and is
maintained for backward compatibility, so that dependency needs its own
qualification \citep{sw_synthdid,sw_gsynth}. HonestDiD takes an estimated
event-study vector and covariance and adds an explicit class of possible
trend violations \citep{sw_honestdid}. That is a sensitivity operation on
an already defined contrast, not an automatic selection of a true trend
restriction.

CausalImpact and CausalPy provide further useful comparisons
\citep{sw_causalimpact,sw_causalpy}. CausalImpact's primary documentation
requires unaffected controls and a stable counterfactual relationship
after the intervention. CausalPy offers distinct quasi-experimental
designs with Bayesian or OLS estimation. The appropriate comparison with
a cohort estimator or factor model asks whether each model answers the
same bank intervention and whether its counterfactual restrictions are
credible. Comparing interval widths without those restrictions would
reward unwarranted certainty.

\section{Longitudinal libraries must preserve the intervention law}
\label{r3:software-longitudinal}

The longitudinal g-formula can be evaluated through fitted transition
distributions, backward outcome regressions or targeted/augmented
constructions. gfoRmula, ltmle, lmtp and zEpid expose different parts of
this space \citep{sw_gformula,sw_ltmle,sw_lmtp,sw_zepid}. They cannot be
made interchangeable by renaming their output \emph{policy effect}.
The treatment law, censoring intervention, endpoint and history are part
of the mathematical input.

gfoRmula's inspected interface takes person-period data, conditional models,
history functions and named interventions. It returns survival risks or
terminal means/probabilities. The number of simulated persons,
\texttt{nsimul}, controls Monte Carlo approximation; the number of
bootstrap samples, \texttt{nsamples}, concerns sampling uncertainty.
Increasing the former cannot compensate for a small observed cohort.
The interface also exposes competing-event treatment, visit processes and
prebaseline lag conventions. Its \texttt{sim\_trunc} option defaults to
true for specified continuous covariate families, restricting simulated
values to the observed range. Such a transformation must be part of the
reported simulation law. For a simple project example,
$P(L=-2,0,2)=(1/4,1/2,1/4)$ gives $\E[L^2]=2$; clipping to $[-1,1]$
gives $1/2$. The numeric range may look safer while the tail-sensitive
financial target changes.

ltmle represents chronology through column order in wide data. Its
treatment nodes must be binary, its survival indicators must remain one
after the event, and its probability bounds concern \emph{cumulative}
treatment/censoring factors. Its output retains influence values and
both bounded and unbounded cumulative probabilities. Continuous outcomes
can be scaled and, when \texttt{Yrange} is supplied, truncated. Some
robust-variance options are unavailable with nonbinary outcomes,
G-computation or other listed options \citep{sw_ltmle}. A sensible adapter
constructs columns from a separately validated clock and records each of
these choices. Renaming an arbitrary real-valued bank action as a binary
node would create a different intervention.

lmtp supplies targeted and sequentially doubly robust estimators for
modified treatment policies, with explicit ordered histories, censoring,
competing events, sampling weights and optional cluster identifiers.
Its finite-history parameter $k$ is a modeling restriction when the full
history is relevant. Its source also resolves a potentially serious
terminological ambiguity. The function \texttt{ipsi} implements a
stochastic retention-or-replacement rule. For
$e(h)=P(A=1\mid H=h)$, direct integration over the auxiliary uniform
random variable in the source gives
\begin{equation}
 q_\delta(1\mid h)=
 \begin{cases}
  \delta e(h),&0\leq\delta\leq1,\\
  1-\{1-e(h)\}/\delta,&\delta>1.
 \end{cases}
 \label{r3:software-ipsi}
\end{equation}
The downward branch scales exposure risk; the upward branch scales
nonexposure risk by $1/\delta$. This differs from the odds multiplier
$\delta e/(1-e+\delta e)$ in Section~\ref{r3:incremental}.
At $e=.2$, $\delta=.5$ produces $.1$ instead of $1/9$;
$\delta=2$ produces $.6$ instead of $1/3$. The corresponding policy means
and influence functions therefore cannot be interchanged. The comparison
is a derivation from inspected source, not a package execution.

The score difference can be made explicit. For a one-time binary action
and a differentiable probability transform $q=f(e)$, let
$m_a(h)=\E[Y\mid A=a,H=h]$ and
$\mu_f(h)=(1-f(e(h)))m_0(h)+f(e(h))m_1(h)$.
Differentiating the outcome regressions, the propensity and the marginal
history law separately, as in Section~\ref{r3:incremental}, gives
\begin{align*}
 \phi_f={}&\frac{Af(e)}{e}(Y-m_1)
       +\frac{(1-A)(1-f(e))}{1-e}(Y-m_0)\\
       &+f'(e)(m_1-m_0)(A-e)+\mu_f-\Psi_f .
\end{align*}
Here histories with $0<e<1$ are considered. For
Equation~\eqref{r3:software-ipsi}, $f'(e)$ is $\delta$ on the downward
branch and $1/\delta$ on the upward branch. For odds tilting it is
$\delta/\{1+(\delta-1)e\}^2$. At $\delta=1$ all reduce to the natural-course
mean and the score reduces to $Y-\E[Y]$. This limiting check and a finite-law
directional derivative are useful acceptance criteria for a policy adapter.

zEpid offers useful Python references for weighting, g-computation,
targeting and transport. Its documentation specifically limits its IPMW
description to monotone missingness and its SNMM g-estimation to a single
time point \citep{sw_zepid}. An umbrella description of \emph{g-methods}
must not be promoted to an implementation of every longitudinal method.
For collections, bankruptcy, settlement and exit should be represented
as the substantive events of Section~\ref{r3:competing-events}; a library's
censoring switch does not decide which event the bank intended to remove.

\section{A fitted state model is not yet a causal planner}
\label{r3:software-states}

Dynamax is a natural numerical comparator for conditional state-space
likelihoods. The inspected HMM abstraction builds initial, transition and
emission components, supplies \texttt{marginal\_log\_prob},
\texttt{filter}, \texttt{smoother}, and an EM E-step
\citep{sw_dynamax}. Its optional external inputs enter through those
components; the standard transition model can be a fixed transition
matrix. Thus an input argument does not prove that a selected subclass
has an action-dependent transition.

More fundamentally, the conditional emission likelihood does not
automatically contain a private action logger. Under the common law of
Section~\ref{r3:state-bridge}, observations of an action can alter the
posterior state through the factor $q_\eta(a_t\mid s_t,h_t)$. If logging is
independent of hidden state conditional on the recorded history, that
factor may cancel from the state update. Otherwise it must be represented
in the observational likelihood. Replacing it by a proposed intervention
then changes the experiment. A state-space adapter must establish this
factorization before fitting; optimizer convergence cannot repair an
omitted likelihood factor.

PyMC and NumPyro can express richer probabilistic models and estimate their
parameters \citep{sw_pymc,sw_numpyro}. They provide numerical inference,
not a general identification theorem. PyMC's inspected example distinguishes
model intervention and observation through \texttt{do} and
\texttt{observe}; their causal interpretation still depends on the supplied
model. NumPyro supports HMC/NUTS, mixed discrete/continuous methods and
variational alternatives. Marginalizing a finite hidden state is often
useful, but it is not the only documented approach to discrete variables.
Comparisons should hold the likelihood and prior fixed and separately
measure integration error, parameter uncertainty and sensitivity to
observationally equivalent causal completions.

POMDPs.jl and pomdp\_py solve another layer of the problem
\citep{sw_pomdps,sw_pomdp_py}. They provide problem representations,
planning and simulation interfaces. The Julia core interface, POMDPTools,
a solver and a Python bridge have separate versions and dependencies.
An adapter must state whether the observation occurs before or after the
transition, what information the policy receives and when reward is paid.
A smoothed state using future bank records cannot be used as an online
decision belief. Solver residuals and rollout error should be evaluated
against the model-conditional value, while the identification and model
error are assessed through Chapter~\ref{r3:controlled-identification}.

d3rlpy, SCOPE-RL and Open Bandit Pipeline offer policy-learning and
evaluation comparators \citep{sw_d3rlpy,sw_scope_rl,sw_obp}. Their natural
objects differ: a fitted policy, a sequential return estimate or return
distribution, and a contextual-bandit reward. A one-period contact reward
cannot represent a multimonth recovery path with carryover. Occupancy-ratio
OPE and trajectory-ratio OPE also require different sufficient-information
and coverage conditions. The rare-path construction in
Section~\ref{r3:ope-hardness} is a relevant stress case even when each
logged action has a positive probability.

\section{Discovery outputs must retain their unresolved information}
\label{r3:software-discovery}

The graph is a scientific output with a type. Tigramite's documentation
distinguishes PCMCI, PCMCI+, LPCMCI and other variants, and separately lists
the restrictions of conditional-independence tests \citep{sw_tigramite}.
Selecting a nonlinear test does not remove a stationarity assumption;
adding a missing-value mask does not identify an informative observation
process. A temporal discovery result must retain lag convention, window,
mask, background knowledge, test and model class.

The inspected causal-learn FCI interface returns a graph marked as a PAG
and edge-property objects. Its matrix distinguishes tails, arrowheads
and circles \citep{sw_causallearn}. Converting that matrix into Boolean
adjacency would erase precisely the information needed for
effect identification. A circle endpoint expresses unresolved ancestral
information; it is not a weakly weighted arrow. Bidirected edges are not
two ordinary causal arrows. The receiving effect operation must reason
over compatible graph completions or use a criterion valid for the
equivalence class, as in Chapter~\ref{r3:discovery}.

NOTEARS and DAGMA supply useful continuous-optimization constructions,
while gCastle collects several algorithm families and Tetrad offers a
separate Java ecosystem \citep{sw_notears,sw_dagma,sw_gcastle,sw_tetrad}.
For the optimization methods, the acyclicity constraint is a statement
about a fitted graph; it is not a proof that the observed law uniquely
identifies that graph. Search tolerances, thresholding and initialization
belong to the numerical output. For a toolbox, the assumptions belong to
the selected algorithm, not to the package name. gCastle's inspected
module distinguishes IID and event-sequence algorithms and marks HPCI
as under development. Tetrad's language bridge and runtime need separate
qualification from a particular search procedure.

Causica combines graph and neural-mechanism fitting within an additive-noise
structural model, then simulates interventions \citep{sw_causica}.
This is a useful comparison where the structural restrictions are justified.
A variational graph distribution conditional on that model is not a
frequentist graph-confidence set or a distribution-free certificate of
causal identification. The current README explicitly says Rhino and
DDECI were removed from the codebase and refers to an earlier release.
Their paper descriptions cannot establish the capability of the current
snapshot. The initial CausalNex repository retrieval failed; that historical
failure is not evidence that the project is abandoned \citep{sw_causalnex}.
The later pinned-source comparison recovers its StructureModel and selected
NOTEARS implementations. Default adoption remains deferred pending a justified
consumer, an isolated compatible environment and dependency qualification;
source retrieval itself is no longer the obstacle
\citep{exec13_pydisc_causalnex}.

\section{Data organization must preserve what was known}
\label{r3:software-data}

DuckDB and Polars provide relational and columnar execution paths;
Pandera supplies explicit schema and value checks
\citep{sw_duckdb,sw_polars,sw_pandera}. These are useful building blocks
for Chapter~\ref{r3:observation}, but type-correct data may still encode
the wrong decision history. An event dated before a credit-limit decision
may have reached the bank after it. A later corrected value can also
replace the value actually used. The history constructor must therefore
enforce event, availability and revision semantics before exporting
predictors to an estimator.

Feast's inspected historical-retrieval API makes the distinction concrete.
It uses event time and TTL, and offers
\texttt{filter\_by\_created\_timestamp}. The latter defaults to false;
when true it excludes later-created values only for feature views with
a configured creation-timestamp field \citep{sw_feast}. Hence it is
incorrect either to assume every point-in-time join captures actual
availability or to claim this capability is entirely absent. The bank
must map its availability semantics to the configured fields and test a
backfill whose event precedes, but creation follows, the decision.
TTL governs freshness and retrieval scope; it cannot establish that mapping.

DataHub and OpenMetadata can supply asset descriptions, ownership,
glossaries and lineage \citep{sw_datahub,sw_openmetadata}. Their graphs
record how data are described and produced; they are not causal graphs
of the customer or economy. A transform that computes arrears from
transactions produces a lineage edge, not evidence that editing the
arrears column would change past transactions. The product should reuse
the bank's catalog and attach the causal interpretation to its stable
asset identifiers. Broad catalog marketing that mixes open-source and
managed-cloud features is not evidence that a particular local deployment
has those features.

OpenLineage supplies an extensible run/job/dataset event model
\citep{sw_openlineage}. A causal-analysis extension should preserve the
target, data snapshot, observation clock, model/graph, split plan and
estimator identity. These are proposed scientific facets, not standard
causal semantics already supplied by OpenLineage. Splink addresses another
part of the observation operator: probabilistic linkage supplies pair
probabilities and estimated entity clusters \citep{sw_splink}. Their
uncertainty matters when a false household merge changes the unit of
treatment or uses post-treatment information. The linkage-variance
decomposition in Section~\ref{r3:entity-uncertainty} explains why a
thresholded ID alone is not a sufficient uncertainty record.

\section{Distributed execution must preserve the score experiment}
\label{r3:software-execution}

After identification and nuisance fitting, some estimates can be reduced
through additive scores. This creates an opportunity for local SQL,
columnar or distributed execution. Ray supplies tasks, actors and shared
objects, while Polars documents streaming and optional GPU execution
\citep{sw_ray,sw_polars}. Neither system specifies the independent
information in a causal analysis.

For a scalar estimating equation $\sum_{i=1}^n\psi_i(\widehat\theta)=0$,
write $\widehat J=n^{-1}\sum_i\partial_\theta\psi_i$.
If sampling clusters $c$ are independent, a basic cluster sandwich uses
\begin{equation}
 \widehat{\operatorname{Var}}(\widehat\theta)=
 \widehat J^{-2}\frac1{n^2}\sum_c
       \left(\sum_{i\in c}\psi_i\right)^2,
 \label{r3:software-cluster-reduction}
\end{equation}
before any stated finite-cluster correction. The reduction key is the
statistical cluster, not the worker partition. With score values
$(1,1,-1,-1)$ in two clusters $(1,1)$ and $(-1,-1)$, the cluster meat is
$8$; treating rows as independent gives $4$. Sharding each cluster across
workers and squaring partial sums before regrouping repeats that error.
Correct distributed execution first sums by stable cluster identity and
then forms the outer products. Multiway clustering and serial dependence
require their own covariance constructions.

Cross-fitting has the same constraint. A fold assignment must be a durable
function of the permitted sampling unit or time block, fixed before
worker scheduling. A retry must not duplicate a score contribution;
repartitioning must not move a household between training and evaluation.
Floating-point reductions may differ across orders, so a declared numerical
tolerance is appropriate, but the statistical partition must be exact.

The resulting qualification sequence is constructive. First verify a
small finite-law functional and its failure cases. Then compare the
adapter's actual target and outputs against a numerical reference.
Next evaluate sampling bias, coverage or policy regret under named
experiments. Finally measure the same scientifically valid computation
under larger loads, including memory, spill, retries and cost. A larger
row count is useful only after these identities are preserved.

\section{What the comparison assigns to the shared product}
\label{r3:software-composition}

The reusable layer should own the intervention and observation
representations, target-preserving transformations, identification evidence,
inference requirements and immutable execution records. Existing libraries
can then supply well-defined numerical operations. A bounded first set
includes adjustment-based effects, valid linear-IV moments, specified
panel/cohort contrasts, finite-horizon observed-history policies and
graph-equivalence discovery outputs. Controlled latent models require a
separate observational/interventional likelihood interface and evidence
for the requested value; supplying a planner alone cannot complete it.

This division also provides a substantive alternative to building the
product. An expert team can compose the same open-source libraries directly.
The shared product has value only if it preserves the scientific contracts
more reliably, reduces repeated data and validation work, and improves
decision throughput at a measured total cost. An estimator wrapper with
a common name adds little. The operations, invalid requests, numerical
obligations and reusable artifacts in Chapter~\ref{r3:product} make this
claim testable. The comparisons here explain why those requirements follow
from the economic question, rather than from a preference for a particular
software interface.

